\bigskip\Needspace{10\baselineskip}
\pdfbookmark[1]{Pallas Olbersiana. November 1802}{pallas180211}
\begin{center}
PALLAS OLBERSIANA.\par
\medskip\hrule\medskip
Monatliche Correspondenz. Vol. VI. pp.\ 499—505. November 1802.\par
\medskip\hrule
\end{center}\medskip
—\quad —\quad —

Dr.\ \textsc{Herschel} was kind enough to send us an offprint of his paper from the \textit{Philosophical Transactions}\HistoricalNote{22}. We already gave our readers its essential contents in an extract from one of his letters in the July issue, p.\ 90. Having read the paper through carefully again, we see no reason to change the opinion we have already expressed on it. We also see from our correspondence that other astronomers, including \textsc{La Lande}, \textsc{Oriani}, \textsc{Piazzi} and \textsc{Gauss}, do not agree with the new classification he proposes. Dr.\ \textsc{Gauss}’s arguments and a particularly apt observation of his against the supposed so-called harmonic law with which the planets’ distances are said to agree deserve special consideration. This perceptive mathematician expresses himself on the matter in a remarkable way in a letter dated 16 October.

“Dr.\ \textsc{Herschel}, as Prof.\ \textsc{Huth}, who visited him in England, told me, still will not tolerate the new planets, although, as far as I know, not a single astronomer has yet approved his proposal. As you so rightly observed, whether we choose to regard \textit{Ceres} and \textit{Pallas} as planets is at bottom a matter of convention. The issue is in no way whether they \textit{are} planets, but whether it is \textit{proper} and \textit{fitting} to \textit{call} these celestial bodies planets. They agree with the \textit{previously known} planets in some respects and differ from them in others. In this year’s July issue you have shown clearly enough that the latter points are quite immaterial. (They cannot even include the inclination of the orbit in the case of \textit{Ceres}, as you observed in last year’s December issue.) What seems to show that astronomers have always regarded a nearly circular orbit and the resulting \textit{enduring presence} as essential is that
\SourceLeaf{236}%
they recognized them as planets at once without further question, as soon as they were satisfied that the orbit had that character. I even think that if the future should confirm the conjecture of our excellent \textsc{Olbers} that \textit{Ceres} and \textit{Pallas} are only fragments of a shattered planet, through the discovery of other fragments, we would still not necessarily have to give up calling them planets. For it seems to me that in examining their claim to the title of planet, what matters more is \textit{whether} these celestial bodies are planets by virtue of their essential property than \textit{how} they \textit{became} planets.

“It is strange that anyone should have wanted to use the so-called law given by Prof.\ \textsc{Titius} as an argument against the two planets. Contrary to the nature of all truths that deserve to be called laws, this relation holds only incidentally for the other planets. For Mercury it does not hold at all, a point that seems to have gone entirely unnoticed so far\HistoricalNote{23}. It seems quite evident to me that the sequence
\[
4,\ 4+3,\ 4+6,\ 4+12,\ 4+24,\ 4+48,\ 4+96,\ 4+192
\]
with which the distances were supposed to agree is not even a continuous sequence. The term preceding $4+3$ cannot be $4$, that is, $4+0$. It has to be $4+1\frac12$. Thus there should be infinitely many terms between $4$ and $4+3$. Or, as \textsc{Wurm} puts it (M.\ C., June 1801 issue, p.\ 594), for Mercury, or for $n=1$, the expression $4+2^{n-2}\,3$ does not give $4$, but $5\frac12$.

“There is nothing to criticize in looking for such approximate correspondences in nature. The greatest men have indulged in such \textit{lusus ingenii}. For all the pride \textsc{Kepler} took in his regular solids, which he had brought into agreement with the planetary distances (he said he would not give up the honor of this discovery for the Electorate of Saxony), he would certainly not have used them to dispute Uranus’s status as a planet if its discovery had been made in his day, merely because Uranus did not fit his ideas. Without any doubt he would instead have abandoned those ideas at once. You did just that. You too considered that law of \textsc{Titius}, but you called your reflections mere dreams and dropped them as soon as they began to contradict the facts.”
\begin{flushright}[\textsc{von Zach.}]\end{flushright}

